% ---------------------------------------------------------------------------
% Resultados y discusión
% ---------------------------------------------------------------------------
% Mínimo dos páginas por capítulo (norma 2.6).

\chapter{Resultados y discusión}
\label{cap:resultados}

Este capítulo presenta los resultados de la campaña de evaluación de la Fase 6, ejecutada el 2 de
septiembre de 2026 sobre la partición de evaluación del conjunto de alertas etiquetado ---205
alertas, de las cuales 18 pertenecen al caso de uso acotado y tienen, por tanto, una clase
verdadera contra la que se puede medir---, con el perfil de cliente empresarial y el justificador de
mil millones de parámetros a temperatura de muestreo nula.

% ===========================================================================
\section{Resultado principal: clasificación frente al método de referencia}
% ===========================================================================

La pregunta que la fase debía responder es si el triaje asistido aporta algo sobre el nivel de
regla de Wazuh tomado como método tradicional de referencia. La tabla \ref{tab:resultado-principal}
resume el resultado.

\begin{table}[H]
  \centering
  \caption{Prototipo frente al baseline de Wazuh (partición de evaluación, $n=205$)}
  \label{tab:resultado-principal}
  \begin{tabular}{@{}lcc@{}}
    \toprule
    \textbf{Métrica} & \textbf{Prototipo} & \textbf{Baseline (nivel $\geq 5$)} \\
    \midrule
    Precisión & 0,556 & 0,154 \\
    Exhaustividad & 1,000 & 1,000 \\
    Medida $F_1$ & 0,714 & 0,267 \\
    Tasa de falsos positivos & 0,041 & 0,282 \\
    \bottomrule
  \end{tabular}
\end{table}

La matriz de confusión del prototipo es $VP=10$, $FP=8$, $VN=187$, $FN=0$; la del baseline en su
umbral óptimo es $VP=10$, $FP=55$, $VN=140$, $FN=0$. Ambos métodos detectan la totalidad de las
amenazas reales de la partición ---exhaustividad de 1,0 en los dos casos---, pero el prototipo reduce
la tasa de falsos positivos de 0,282 a 0,041, casi siete veces menos. La diferencia son 187 eventos
de ruido de plataforma que el motor clasifica correctamente como no soportados y que el nivel de
regla, por sí solo, no tiene forma de suprimir. Esta reducción del ruido operativo sin pérdida de
exhaustividad es, precisamente, el resultado que el proyecto se propuso demostrar.

% ===========================================================================
\section{El límite de precisión del prototipo, y por qué}
% ===========================================================================

La precisión del prototipo se queda en 0,556, no en 1,0, por una razón concreta y verificable: las
dieciocho alertas soportadas de la partición corresponden todas al mismo activo y al mismo
servicio ---acceso por SSH al servidor con vulnerabilidades documentadas del laboratorio---, de las
cuales diez son intentos de acceso reales y ocho son actividad de un administrador legítimo que se
equivocó de credencial. El clasificador determinista decide por familia de alerta y por postura de
exposición del activo, y como el servicio está efectivamente expuesto en ese nodo, marca las
dieciocho como amenaza potencial: no dispone de ninguna señal sobre la legitimidad del origen que le
permita distinguir al administrador del atacante. Es el límite esperado de un clasificador que no
incorpora aprendizaje sobre patrones de comportamiento, y es exactamente la limitación que
motivaría un clasificador con ajuste fino si el conjunto de datos tuviera variedad suficiente de
familias para entrenarlo sin invalidar la propia evaluación.

% ===========================================================================
\section{Por qué el óptimo del baseline es débil}
% ===========================================================================

El nivel de regla de Wazuh se evaluó como clasificador barriendo su umbral de decisión; la tabla
\ref{tab:barrido-baseline} resume la medida $F_1$ resultante en cada punto.

\begin{table}[H]
  \centering
  \caption{Medida $F_1$ del baseline según el umbral de nivel de regla}
  \label{tab:barrido-baseline}
  \begin{tabular}{@{}lcccccccccc@{}}
    \toprule
    nivel $\geq$ & 3 & 4 & 5 & 6 & 7 & 8 & 9 & 10 & 11 & 12 \\
    \midrule
    $F_1$ & 0,093 & 0,267 & 0,267 & 0,035 & 0,035 & 0,167 & 0,167 & 0,167 & --- & --- \\
    \bottomrule
  \end{tabular}
\end{table}

El mejor punto del baseline ---$F_1=0,267$, en los umbrales 4 y 5--- no se acerca en ningún caso al
resultado del prototipo ($F_1=0,714$). La curva, además, no es monótona: subir el umbral no mejora
de forma limpia, porque los verdaderos positivos y buena parte del ruido conviven en los mismos
niveles de severidad. Esto confirma que el nivel de regla, por sí solo, no basta como criterio de
prioridad, y justifica la reclasificación que realiza el motor de triaje.

% ===========================================================================
\section{Operación y continuidad del servicio}
% ===========================================================================

El tiempo de clasificar y decidir sobre una alerta, con el clasificador determinista, fue del orden
de $0,01\,\text{ms}$ por alerta, muy por debajo del presupuesto de latencia establecido. Justificar
con el modelo de lenguaje de mil millones de parámetros tomó, en esta campaña, aproximadamente
$36\,\text{s}$ por alerta ---más que en una invocación aislada, porque cada justificación recarga el
modelo de $808\,\text{MB}$ en un proceso nuevo; en un despliegue real el modelo permanecería
cargado entre invocaciones---. La cobertura del clasificador ---la proporción de alertas que caen dentro
del caso de uso acotado--- fue de 0,088: solo el 8,8\,\% de las alertas de un Wazuh recién desplegado
corresponden a un patrón clasificable, y el resto es telemetría interna de la propia plataforma. La
tasa de escalado a validación humana fue de 0,0 sobre este conjunto, porque el clasificador
determinista nunca reportó una confianza por debajo del umbral de escalado en estas alertas.

Sobre la métrica de continuidad ---el requisito que distingue a este trabajo de un sistema de
respuesta automatizada convencional---, no se registró ninguna acción disruptiva indebida: ninguna
acción propuesta a partir de una alerta que resultó ser un falso positivo alcanzaba un servicio que
el cliente prestara. El escenario más exigente de esta métrica ---una acción de alto impacto que deba
ser efectivamente retenida por la validación humana--- no llegó a ejercitarse sobre este conjunto de
datos, lo que se documenta como una limitación de la muestra y no como una garantía general.

% ===========================================================================
\section{Calidad de la justificación: anclada, pero inicialmente poco fiable}
% ===========================================================================

La verificación automática de anclaje ---que ninguna dirección IP mencionada sea distinta de la de la
alerta y que el texto referencie al menos un campo concreto de ella--- dio una tasa de 18 sobre 18: el
cien por ciento de las justificaciones del modelo de mil millones de parámetros pasó la
verificación. Sin embargo, la revisión manual de esa misma muestra reveló que un anclaje perfecto no
garantiza una explicación correcta. Un ejemplo textual generado por el modelo antes de incorporar
recuperación aumentada:

\begin{quote}
\itshape
«La regla 5760 es un protocolo de seguridad para la autenticación de usuarios en redes. El MITRE
T1110.001 es un conjunto de técnicas de seguridad para la autenticación de usuarios...»
\end{quote}

La regla 5760 de Wazuh no es un protocolo: es la regla que detecta un fallo de autenticación por
SSH. La técnica T1110.001 no es un conjunto de técnicas de seguridad: es \emph{Password Guessing},
una técnica de ataque del catálogo MITRE ATT\&CK. El modelo cita los campos correctos ---de ahí que
ancle--- pero describe su significado al revés, y de las dieciocho justificaciones generadas solo
cuatro textos resultaron distintos entre sí, porque una temperatura de muestreo nula sobre un único
nodo produce salidas casi idénticas. La conclusión de esta medición es metodológicamente importante:
la tasa de anclaje mide que el modelo no inventa campos, no que acierte su significado, y no debe
leerse como un indicador de calidad de la justificación.

% ===========================================================================
\section{El contraste con recuperación aumentada}
% ===========================================================================

Sobre las mismas dieciocho alertas, con el mismo modelo y la misma temperatura, se repitió la
campaña con el módulo de recuperación aumentada activado. La tasa de anclaje automática no cambió
---se mantuvo en 18 sobre 18, porque mide otra cosa---, pero la corrección semántica de la
justificación, evaluada manualmente, sí lo hizo. La misma alerta que antes producía el texto citado
en la sección anterior, generó con recuperación aumentada:

\begin{quote}
\itshape
«La regla 5760 de Wazuh indica un ataque de fuerza bruta SSH en curso desde una IP.»
\end{quote}

correcta y respaldada por los pasajes recuperados del corpus curado sobre las reglas de fuerza
bruta SSH. En una alerta de falso positivo, el contraste fue igual de marcado: de «la regla 5710
define el protocolo SSH como un protocolo seguro» ---una afirmación sin sentido--- a «la regla 5710 de
Wazuh indica que se probó un nombre de usuario que no existe por SSH, un indicio de ataque de
fuerza bruta», que es la lectura correcta del evento.

Tres matices, medidos y no solo observados, acompañan esta mejora. Primero, la diversidad de
textos no aumentó ---siguieron saliendo solo cuatro textos distintos de dieciocho---, porque la causa de
esa baja diversidad es el propio conjunto de datos concentrado en un nodo, no el mecanismo de
justificación. Segundo, la recuperación resultó imprecisa en el identificador exacto de la regla:
para una alerta de la regla 5760, el sistema recuperó el grupo completo de reglas de fuerza bruta
SSH (5710, 5712, 5763) en lugar del identificador preciso; aun así, el dominio recuperado era el
correcto, y por eso la justificación mejoró. Tercero, con recuperación aumentada el modelo tiende a
parafrasear los pasajes recuperados ---enumerar las reglas del corpus--- más que a razonar de forma
original sobre la alerta concreta, lo que es preferible a inventar una explicación pero sigue
estando lejos de la redacción de un analista humano.

% ===========================================================================
\section{Verificación sobre el conjunto completo}
% ===========================================================================

Para descartar que el resultado dependiera de la partición elegida, la clasificación ---sin el
justificador, por su coste de tiempo--- se repitió sobre las 410 alertas del conjunto completo (36
soportadas). Las tasas obtenidas ---precisión 0,556, exhaustividad 1,0, $F_1$ 0,714, tasa de falsos
positivos 0,041--- resultaron idénticas a las de la partición de evaluación, lo que era el resultado
esperado: el clasificador es determinista y no se entrena, así que no existe diferencia entre
entrenar y evaluar sobre él. La partición se respetó de todas formas por rigor metodológico.

% ===========================================================================
\section{Limitaciones y análisis de errores}
% ===========================================================================

Cuatro limitaciones, todas medidas y no solo declaradas, acotan el alcance de estos resultados.
Primero, el tamaño de la muestra soportada ---dieciocho alertas, de una sola familia de ataque sobre
un único activo--- hace que los porcentajes de precisión se sostengan sobre pocos casos, y que la
correlación de orden de prioridad no pudiera calcularse por falta de variedad. Segundo, el error
sistemático dominante del prototipo ---los ocho falsos positivos--- tiene una causa identificada y no
un origen difuso: la ausencia de una señal sobre la legitimidad del origen de la conexión. Tercero,
el clasificador con ajuste fino permanece bloqueado por la misma concentración del conjunto de
datos, así que el desajuste esperado entre un modelo preentrenado sobre prosa y una entrada de
eventos estructurados queda documentado, pero no cuantificado. Cuarto, el auditor de
vulnerabilidades no es infalible: la ausencia de un hallazgo no demuestra que un activo sea seguro,
y la postura de exposición usada como contexto de decisión se tomó de los hallazgos reales de una
única campaña de escaneo.

El resultado global, no obstante, es positivo y medido sobre el problema que el proyecto se propuso
atacar: el triaje asistido reduce el ruido operativo de forma sustancial sin sacrificar
exhaustividad y sin comprometer la continuidad del servicio, y su límite de precisión y su
justificador tienen, ambos, una causa identificada y un camino de mejora conocido ---un clasificador
entrenado con datos de mayor variedad y un modelo de justificación de mayor capacidad---, que la
propia interfaz de análisis del prototipo ya permite incorporar sin rediseñar el sistema.
